\documentclass[12pt,a4paper]{article}

\usepackage[T2A]{fontenc}
\usepackage[utf8]{inputenc}
\usepackage[russian]{babel}
\usepackage{amsmath,amssymb,amsthm,mathtools}
\usepackage{geometry}
\usepackage{microtype}
\usepackage{enumitem}
\usepackage{xcolor}
\usepackage{tikz}
\usepackage{hyperref}

\geometry{left=25mm,right=20mm,top=22mm,bottom=24mm}
\hypersetup{
  colorlinks=true,
  linkcolor=blue!45!black,
  urlcolor=blue!55!black
}
\setlist{nosep}
\setlength{\parindent}{1.25em}
\setlength{\parskip}{0.15em}

\newtheorem{theorem}{Теорема}
\newtheorem{proposition}{Утверждение}
\newtheorem{lemma}{Лемма}
\theoremstyle{definition}
\newtheorem{definition}{Определение}
\theoremstyle{remark}
\newtheorem*{remark}{Замечание}

\DeclareMathOperator{\Sp}{Sp}
\DeclareMathOperator{\diag}{diag}
\DeclareMathOperator{\cond}{cond}
\newcommand{\CC}{\mathbb{C}}
\newcommand{\RR}{\mathbb{R}}
\newcommand{\ZZ}{\mathbb{Z}}
\newcommand{\eps}{\varepsilon}
\newcommand{\norm}[1]{\left\lVert #1\right\rVert}

\DeclareMathOperator{\ind}{ind}
\newcommand{\Conv}{\mathcal C}

\title{Факторизация в алгебре Винера\\Операторы свёртки и проекторы}
\author{Конспект лекции}
\date{24 октября 2026 года}

\begin{document}
\maketitle

\section{Факторизация функции, близкой к единице}

Напомним обозначения предыдущей лекции. Алгебра Винера состоит
из функций на $\Gamma=\{t\in\CC:|t|=1\}$ вида
\[
  a(t)=\sum_{k\in\ZZ}a_kt^k,\qquad
  \norm{a}_W=\sum_{k\in\ZZ}|a_k|<\infty.
\]
Подалгебры $W_+$ и $W_-$ содержат соответственно неотрицательные
и неположительные степени $t$. Обе замкнуты в $W$ и содержат единицу.

\begin{theorem}
Если $d\in W$ и $\norm{d}_W<1$, то
\[
  \boxed{1-d(t)=c_-(t)c_+(t)},\qquad
  c_\pm,c_\pm^{-1}\in W_\pm.
\]
При этом $\ind(1-d)=0$.
\end{theorem}

Геометрическая часть уже обсуждалась: $|d(t)|\leqslant\norm{d}_W<1$,
поэтому образ $1-d$ лежит в правой полуплоскости. Его непрерывный
аргумент после обхода возвращается к начальному значению.
Теперь построим сами множители.

\subsection{Логарифмический ряд в норме Винера}

Скалярные разложения имеют вид
\[
  \log(1-z)=-\sum_{n=1}^{\infty}\frac{z^n}{n},\qquad
  \log(1+z)=\sum_{n=1}^{\infty}\frac{(-1)^{n+1}z^n}{n},
  \qquad |z|<1.
\]
Их радиус сходимости равен $1$. Подстановка элемента $d$ корректна,
поскольку ряд сходится абсолютно в банаховой алгебре:
\[
  h:=\log(1-d)=-\sum_{n=1}^{\infty}\frac{d^n}{n}\in W,
  \qquad
  \norm{h}_W\leqslant\sum_{n=1}^{\infty}\frac{\norm{d}_W^n}{n}<\infty.
\]
В каждой точке окружности это логарифм, заданный указанным степенным
рядом; поэтому $e^{h(t)}=1-d(t)$.

Разложим $h$ в абсолютно сходящийся ряд Фурье:
\[
  h(t)=\sum_{k\in\ZZ}h_kt^k=h_-(t)+h_+(t),
\]
где выберем
\[
  h_-(t)=\sum_{k<0}h_kt^k,\qquad
  h_+(t)=\sum_{k\geqslant0}h_kt^k.
\]
Тогда $h_-\in W_-$ и $h_+\in W_+$. В рукописи постоянный коэффициент
включён в отрицательную часть; здесь он включён в положительную,
чтобы сохранить нормировку предыдущего конспекта.

\subsection{Экспоненты и обратные множители}

Экспоненциальный ряд
\[
  e^u=\sum_{n=0}^{\infty}\frac{u^n}{n!}
\]
сходится в любой банаховой алгебре с единицей для каждого $u$.
При нормировке $\norm{1}=1$, выполненной в $W$ и $W_\pm$, имеем
\[
  \sum_{n=0}^{\infty}\left\lVert\frac{u^n}{n!}\right\rVert
  \leqslant e^{\norm{u}}.
\]
Следовательно, можно положить
\[
  c_-=e^{h_-},\qquad c_+=e^{h_+}.
\]
Замкнутость подалгебр даёт
\[
  c_\pm\in W_\pm,\qquad c_\pm^{-1}=e^{-h_\pm}\in W_\pm.
\]
Поточечное умножение коммутативно, поэтому
\[
  c_-c_+=e^{h_-+h_+}=e^h=1-d.
\]
Это завершает доказательство. При выбранном распределении
постоянного коэффициента $h_-(\infty)=0$, а значит, $c_-(\infty)=1$.

\begin{remark}
Здесь требовалось $\norm{d}_W<1$, а не только $|d(t)|<1$ на окружности:
оценка в норме алгебры обеспечивает сходимость логарифмического ряда
именно в $W$. В формуле обратного множителя существенен знак минус
в экспоненте.
\end{remark}

\section{Общая теорема о факторизации}

\begin{theorem}
Если $a\in W$ и $a(t)\ne0$ при $|t|=1$, то существуют функции
$a_-,a_+$ такие, что
\[
  \boxed{a(t)=a_-(t)t^{\kappa}a_+(t)},\qquad
  a_\pm,a_\pm^{-1}\in W_\pm,\qquad \kappa=\ind a.
\]
\end{theorem}

Доказательство сводит произвольную функцию к произведению
рациональной функции и функции, близкой к единице. Используются
теорема Винера, устойчивость обратимости при малом возмущении
и рациональная факторизация из предыдущей лекции.

\subsection{Приближение конечной суммой Лорана}

По теореме Винера $a^{-1}\in W$. Выберем число
\[
  0<q<\frac12.
\]
Поскольку $\sum_k|a_k|<\infty$, можно найти $N$ так, чтобы
\[
  \sum_{|k|>N}|a_k|<\frac{q}{\norm{a^{-1}}_W}.
\]
Обозначим
\[
  a_N(t)=\sum_{k=-N}^{N}a_kt^k,\qquad
  \Delta(t)=\sum_{|k|>N}a_kt^k.
\]
Тогда $a=a_N+\Delta$ и
\[
  \norm{a-a_N}_W=\norm{\Delta}_W
  <\frac{q}{\norm{a^{-1}}_W}.
\]
Конечная сумма Лорана $a_N$ является рациональной функцией:
$a_N(t)=t^{-N}p_N(t)$, где $p_N$ --- многочлен.

По теореме об устойчивости обратимости элемент $a_N$ обратим в $W$,
причём
\[
  \norm{a_N^{-1}}_W\leqslant\frac{\norm{a^{-1}}_W}{1-q}.
\]
В частности, $a_N$ не обращается в нуль на окружности, поэтому к нему
применима теорема о рациональной факторизации.

\subsection{Выделение малого множителя}

Запишем
\[
  a=a_N(1+a_N^{-1}\Delta)=a_N(1-d),\qquad
  d=-a_N^{-1}\Delta.
\]
По субмультипликативности нормы
\begin{align*}
  \norm{d}_W
  &\leqslant\norm{a_N^{-1}}_W\norm{\Delta}_W\\
  &<\frac{\norm{a^{-1}}_W}{1-q}
       \frac{q}{\norm{a^{-1}}_W}
   =\frac{q}{1-q}<1.
\end{align*}
Именно ради последнего неравенства было выбрано $q<1/2$.
По доказанной выше теореме
\[
  1-d=c_-c_+,\qquad c_\pm,c_\pm^{-1}\in W_\pm,
  \qquad \ind(1-d)=0.
\]
Аддитивность индекса даёт
\[
  \ind a=\ind a_N+\ind(1-d)=\ind a_N.
\]

\subsection{Сборка факторизации}

Рациональную функцию $a_N$ факторизуем:
\[
  a_N=b_-t^\kappa b_+,\qquad
  b_\pm,b_\pm^{-1}\in W_\pm,\qquad
  \kappa=\ind a_N=\ind a.
\]
Тогда
\begin{align*}
  a&=(b_-t^\kappa b_+)(c_-c_+)\\
   &=(b_-c_-)t^\kappa(b_+c_+).
\end{align*}
Положим $a_-=b_-c_-$, $a_+=b_+c_+$. Произведения и их обратные
принадлежат соответствующим подалгебрам:
\[
  a_\pm\in W_\pm,\qquad
  a_\pm^{-1}=c_\pm^{-1}b_\pm^{-1}\in W_\pm.
\]
Теорема доказана.

\begin{remark}
В доказательстве существенно приближение в норме Винера. Одного
равномерного приближения функции на окружности недостаточно для
приведённой оценки $\norm{d}_W$. Теорема Винера здесь используется
как ранее сформулированный результат, без повторного доказательства.
\end{remark}

\section{Обратимость двустороннего оператора свёртки}

Пусть $a=(a_k)_{k\in\ZZ}\in\ell^1(\ZZ)$. Рассмотрим ограниченный
оператор на $\ell^1(\ZZ)$:
\[
  Ax=a*x,\qquad (Ax)_k=\sum_{j\in\ZZ}a_{k-j}x_j.
\]
Его символ --- функция
\[
  a(t)=\sum_{k\in\ZZ}a_kt^k\in W.
\]
Матрица имеет постоянные диагонали: её элемент в строке $k$
и столбце $j$ равен $a_{k-j}$. Центральный фрагмент выглядит так:
\[
  A=\begin{pmatrix}
    \ddots&\vdots&\vdots&\vdots&\cdots\\
    \cdots&a_0&a_{-1}&a_{-2}&\cdots\\
    \cdots&a_1&a_0&a_{-1}&\cdots\\
    \cdots&a_2&a_1&a_0&\cdots\\
    \cdots&\vdots&\vdots&\vdots&\ddots
  \end{pmatrix}.
\]
Это двусторонняя бесконечная матрица: индексы строк и столбцов
пробегают все целые числа.

\begin{theorem}
Оператор $A$ обратим на $\ell^1(\ZZ)$ тогда и только тогда, когда
его символ нигде не равен нулю на $\Gamma$.
Если это условие выполнено, то $A^{-1}$ также является оператором
свёртки, а его символ равен $1/a(t)$.
\end{theorem}

\begin{proof}
\emph{Достаточность.} По теореме Винера
\[
  \frac1{a(t)}=b(t)=\sum_{k\in\ZZ}b_kt^k,\qquad (b_k)\in\ell^1(\ZZ).
\]
Пусть $Bx=b*x$. Символ композиции равен произведению символов,
поэтому $AB=BA=I$: коэффициенты функции $a(t)b(t)=1$ образуют
единичное ядро. Следовательно, $B=A^{-1}$.

\emph{Необходимость.} Этот шаг дополняет формулировку рукописи.
Если $a(t_0)=0$ при $|t_0|=1$, рассмотрим ограниченный ненулевой
функционал
\[
  \chi_{t_0}(x)=\sum_{j\in\ZZ}x_jt_0^j.
\]
Абсолютная сходимость двойной суммы позволяет переставить слагаемые:
\[
  \chi_{t_0}(Ax)=a(t_0)\chi_{t_0}(x)=0.
\]
Значит, образ $A$ лежит в ядре $\chi_{t_0}$. Но для последовательности
$e^{(0)}$, равной $1$ в нуле и $0$ в остальных координатах,
$\chi_{t_0}(e^{(0)})=1$. Уравнение $Ax=e^{(0)}$ неразрешимо,
то есть $A$ не сюръективен и потому необратим.
\end{proof}

Если условие теоремы выполнено, решение $Ax=y$ задаётся формулой
\[
  x=A^{-1}y=b*y,\qquad x_k=\sum_{j\in\ZZ}b_{k-j}y_j.
\]

\section{Факторизация оператора свёртки}

Определим двусторонний сдвиг $V$ на $\ell^1(\ZZ)$:
\[
  (Vx)_k=x_{k-1}.
\]
Это свёртка с ядром $e^{(1)}$, поэтому её символ равен $t$.
Сдвиг обратим: $(V^{-1}x)_k=x_{k+1}$. Для любого целого $m$
\[
  (V^mx)_k=x_{k-m},\qquad \norm{V^m}=1.
\]

\begin{theorem}
Если символ оператора свёртки $A$ не обращается в нуль на окружности,
то
\[
  \boxed{A=A_-V^\kappa A_+},\qquad \kappa=\ind a,
\]
где $A_\pm$ и их обратные --- операторы свёртки с символами
из $W_\pm$.
\end{theorem}

\begin{proof}
Факторизуем символ по общей теореме:
\[
  a(t)=a_-(t)t^\kappa a_+(t),\qquad a_\pm,a_\pm^{-1}\in W_\pm.
\]
Пусть $A_\pm$ --- операторы с символами $a_\pm$. Соответствие
между символами и операторами сохраняет умножение, поэтому
$A=A_-V^\kappa A_+$. Символы $a_\pm^{-1}$ задают обратные операторы.
\end{proof}

В матричной записи $A_+$ имеет нулевые элементы выше главной диагонали,
а $A_-$ --- ниже неё. Например, для $A_+$ при $k<j$ индекс $k-j<0$,
и соответствующий коэффициент символа равен нулю.
Такие же ограничения на диагонали выполнены у обратных множителей.

Из факторизации следует дополнительная формула решения:
\[
  A^{-1}=A_+^{-1}V^{-\kappa}A_-^{-1},\qquad
  x=A_+^{-1}V^{-\kappa}A_-^{-1}y.
\]
Для двусторонней свёртки ненулевой индекс не мешает обратимости:
сдвиг $V^\kappa$ обратим при любом целом $\kappa$.

\begin{remark}
С этого места в рукописи начинается тема <<Операторы Винера--Хопфа>>.
В предоставленных страницах вводятся необходимые проекторы;
сам оператор Винера--Хопфа ещё не определён. Утверждения об обратимости
выше относятся к двусторонним операторам на $\ell^1(\ZZ)$.
\end{remark}

\section{Проекторы: определение и основные свойства}

Пусть $X$ --- нормированное пространство, $I$ --- тождественный
оператор на $X$.

\begin{definition}
Ограниченный линейный оператор $P\colon X\to X$ называется
\emph{проектором}, если
\[
  \boxed{P^2=P}.
\]
\end{definition}

Равенство означает, что повторное применение $P$ ничего не меняет.
Если $y=Px$, то $Py=P^2x=Px=y$: на своём образе проектор действует
тождественно. В общем случае речь не идёт об ортогональной проекции.

\begin{figure}[ht]
  \centering
  \begin{tikzpicture}[>=stealth,scale=1.45]
    \draw[->,gray!65] (-0.4,0) -- (2.7,0);
    \draw[->,gray!65] (0,-0.4) -- (0,2.65);
    \draw[thick,blue!60!black] (-0.5,-0.175) -- (2.7,0.945)
       node[right] {$\operatorname{Im}P$};
    \draw[thick] (-0.15,-0.5) -- (0.78,2.6)
       node[above] {$\ker P$};
    \coordinate (p) at (1.5,0.525);
    \coordinate (q) at (0.45,1.5);
    \coordinate (x) at (1.95,2.025);
    \draw[dashed] (p) -- (x) -- (q);
    \draw[->,blue!60!black] (0,0) -- (p);
    \draw[->] (0,0) -- (q);
    \fill (p) circle (0.035) node[below right] {$Px$};
    \fill (q) circle (0.035) node[left] {$(I-P)x$};
    \fill (x) circle (0.035) node[right] {$x$};
    \node[below left] at (0,0) {$0$};
  \end{tikzpicture}
  \caption{Проекция на образ вдоль ядра: $x=Px+(I-P)x$}
\end{figure}

\begin{proposition}
Тождественный оператор $I$ и нулевой оператор $0$ являются проекторами.
Если $P$ --- проектор, то $Q=I-P$ также является проектором.
\end{proposition}

\begin{proof}
Первые два утверждения следуют из $I^2=I$ и $0^2=0$. Далее
\[
  (I-P)^2=I-2P+P^2=I-P.
\]
Кроме того,
\[
  P(I-P)=(I-P)P=0.
\]
\end{proof}

\begin{proposition}
Если $P\ne0$, то $\norm{P}\geqslant1$.
\end{proposition}

\begin{proof}
Из $P^2=P$ и субмультипликативности операторной нормы следует
\[
  \norm{P}=\norm{P^2}\leqslant\norm{P}^2.
\]
Так как $\norm{P}>0$, можно разделить на $\norm{P}$.
\end{proof}

\section{Разложение пространства по проектору}

\begin{theorem}
Для любого проектора $P\colon X\to X$ выполнено
\[
  \boxed{X=\operatorname{Im}P\oplus\ker P}.
\]
Каждый $x\in X$ единственным образом представляется в виде
\[
  x=x_0+x_1,\qquad x_0=Px\in\operatorname{Im}P,\qquad
  x_1=(I-P)x\in\ker P.
\]
\end{theorem}

\begin{proof}
Положим $x_0=Px$ и $x_1=x-Px$. Тогда
\[
  Px_1=P(x-Px)=Px-P^2x=0,
\]
так что $x_1\in\ker P$. Разложение существует.

Чтобы доказать его единственность, проверим нулевое пересечение.
Пусть $z\in\operatorname{Im}P\cap\ker P$. Тогда $z=Py$ для
некоторого $y$, но $Pz=0$. Поэтому
\[
  0=Pz=P^2y=Py=z.
\]
Если есть два разложения, разность их первых компонент равна
минус разности вторых и лежит в этом пересечении. Значит,
соответствующие компоненты совпадают.
\end{proof}

\subsection{Образ дополнительного проектора}

В рукописи предложено упражнение:
\[
  \boxed{\ker P=\operatorname{Im}(I-P)}.
\]
Приведём решение как дополнение. Если $z=(I-P)y$, то
$Pz=Py-P^2y=0$, значит, $z\in\ker P$. Обратно, если $Pz=0$,
то $z=(I-P)z$, то есть $z\in\operatorname{Im}(I-P)$.
Оба включения доказаны.

Применяя это равенство к проектору $I-P$, получаем также
\[
  \operatorname{Im}P=\ker(I-P).
\]
Поэтому образ и ядро ограниченного проектора замкнуты: оба являются
ядрами непрерывных операторов. В банаховом пространстве разложение
выше является прямой суммой замкнутых подпространств с ограниченными
координатными проекциями $P$ и $I-P$.

\section{Координатный проектор на положительную полуось}

На $\ell^1(\ZZ)$ определим оператор
\[
  \boxed{(P_+x)_k=
    \begin{cases}
      x_k,&k\geqslant1,\\
      0,&k\leqslant0.
    \end{cases}}
\]
Это соглашение соответствует рукописи: нулевая координата
в положительную часть не входит.

Оператор линеен, а повторное обнуление тех же координат не меняет
результата, поэтому $P_+^2=P_+$. Ограниченность следует из оценки
\[
  \norm{P_+x}_1=\sum_{k=1}^{\infty}|x_k|
  \leqslant\sum_{k\in\ZZ}|x_k|=\norm{x}_1.
\]
Следовательно, $\norm{P_+}\leqslant1$. Проектор ненулевой,
поэтому $\norm{P_+}\geqslant1$ и
\[
  \boxed{\norm{P_+}=1}.
\]
Можно получить нижнюю оценку и непосредственно: на последовательности
$e^{(1)}$ достигается $\norm{P_+e^{(1)}}_1=\norm{e^{(1)}}_1=1$.

Дополнительный проектор $P_-=I-P_+$ оставляет координаты с $k\leqslant0$.
В обозначениях
\begin{align*}
  X_+&=\{x\in\ell^1(\ZZ):x_k=0\text{ при }k\leqslant0\},\\
  X_-&=\{x\in\ell^1(\ZZ):x_k=0\text{ при }k\geqslant1\}
\end{align*}
имеем
\[
  \ell^1(\ZZ)=X_+\oplus X_-,\qquad
  \operatorname{Im}P_+=X_+,\quad\ker P_+=X_-.
\]
Здесь верна и точная аддитивность нормы:
\[
  \norm{x}_1=\norm{P_+x}_1+\norm{P_-x}_1.
\]
Описание $P_-$ и подпространств разворачивает применение общей
теоремы к данному примеру.

\begin{remark}
Не следует смешивать соглашения для $P_+$ и $W_+$.
Подалгебра $W_+$ содержит постоянный коэффициент, тогда как
координатный проектор $P_+$ здесь оставляет только строго
положительные индексы. Ряд, соответствующий $X_+$, не имеет
постоянного члена и принадлежит $tW_+$.
\end{remark}

В конце рукописи ещё раз отмечен общий принцип оценки операторной
нормы: если $\norm{Tx}\leqslant C\norm{x}$ для всех $x$, то
$\norm{T}\leqslant C$.

\end{document}
